\subsection{A short proof of the graded lower comparison}
\label{subsec:selmer-koszul-comparison}

Let $(B,\mathfrak m)$ be a Noetherian local ring, and let
$R=\bigoplus_{\boldsymbol n\ge0}R_{\boldsymbol n}$ be a standard finitely
generated multigraded $B$-algebra, generated in the coordinate degrees $e_i$.
Put $I_{\mathrm{irr}}=\prod_i(R_{e_i})R$.
A homogeneous sequence $\boldsymbol f$, of nonnegative multidegrees,
is admissible on a finite graded module $M$ if
\begin{equation}
 (\mathfrak m R I_{\mathrm{irr}})^N
 \subset\operatorname{Ann}_R M+(\boldsymbol f)
 \quad\text{for some }N.
 \label{eq:selmer-koszul-admissible}
\end{equation}
Its Koszul homology then has finite $B$-length in sufficiently large
degrees. Indeed, the displayed power annihilates that homology; by
standardness, every sufficiently large piece of a finite graded module
is generated by $I_{\mathrm{irr}}^N$. The pieces are therefore killed
by $\mathfrak m^N$. Define
\[
 \chi(\boldsymbol f;M)_{\boldsymbol n}
 =\sum_i(-1)^i\ell_B H_i(K_\bullet(\boldsymbol f;M))_{\boldsymbol n}.
\]
The definition is additive in graded exact sequences and invariant under
permutation of the sections. With $M(-\delta)_{\boldsymbol n}
=M_{\boldsymbol n-\delta}$, the first Koszul factor gives
\begin{equation}
 \chi(f,\boldsymbol t;M)
 =\chi(\boldsymbol t;M/fM)
  -\chi(\boldsymbol t;M[f](-\deg f)).
 \label{eq:selmer-koszul-recursion}
\end{equation}
All lengths in this formula are finite in a common eventual range,
because the tail is admissible on both displayed modules.

\begin{proposition}[Eventual positivity and quotient comparisons]
\label{prop:selmer-koszul-positive}
For an admissible sequence, $\chi(\boldsymbol f;M)\ge0$ in all sufficiently
large degrees. Consequently a graded surjection $M\twoheadrightarrow N$
gives $\chi(\boldsymbol f;M)\ge\chi(\boldsymbol f;N)$.
For an admissible concatenation $(\boldsymbol f,\boldsymbol t)$,
\[
 \chi(\boldsymbol f,\boldsymbol t;M)
 \le\chi(\boldsymbol t;M/(\boldsymbol f)M)
\]
eventually.
\end{proposition}

\begin{proof}
Induct on the number of sections; the empty sequence is a length.
A finite graded prime filtration, additivity, and shifts reduce each
induction stage to $M=R/\mathfrak p$ for a homogeneous prime.
Such a filtration follows by choosing a homogeneous element with maximal
annihilator and applying Noetherian induction. Its annihilator is prime
on homogeneous elements, and hence prime since $\mathbb Z^k$ is ordered.
If $I_{\mathrm{irr}}\subset\mathfrak p$, all sufficiently large pieces
vanish. Otherwise each $R_{e_i}$ contains an element outside
$\mathfrak p$; products give a nonzero element in every nonnegative degree.

Write the first section as $f$ and the tail as $\boldsymbol t$.
If $f\notin\mathfrak p$, multiplication by $f$ on this domain is
injective, and \eqref{eq:selmer-koszul-recursion} invokes the induction
hypothesis on $M/fM$ with one fewer section. If $f\in\mathfrak p$, the
tail is already admissible on $M$. Choose a nonzero $g$ of the same
degree $\delta$. Additivity for
$0\to M(-\delta)\xrightarrow{g}M\to M/gM\to0$ gives
\[
 \chi(f,\boldsymbol t;M)_{\boldsymbol n}
 =\chi(\boldsymbol t;M)_{\boldsymbol n}
  -\chi(\boldsymbol t;M)_{\boldsymbol n-\delta}
 =\chi(\boldsymbol t;M/gM)_{\boldsymbol n}\ge0.
\]
This replacement uses finite-length tail homology; no invariance under
arbitrary changes of a section is asserted. The filtration and induction
use only finitely many bounds, whose componentwise maximum is a common
eventual range. The surjection assertion follows by applying positivity
to its kernel. The last assertion follows by successively using
\eqref{eq:selmer-koszul-recursion} and discarding nonnegative kernel
Euler characteristics.
\end{proof}

\begin{proposition}[Exact scalar extraction]
\label{prop:selmer-koszul-scalar}
Let a degree-zero element $\pi\in B$ act injectively on a finite graded
module $C$, and suppose every graded piece of $C/\pi C$ has finite length.
Set $h(\boldsymbol n)=\ell_B(C/\pi C)_{\boldsymbol n}$. Let $q\ge1$,
let $y_1,\ldots,y_q$ have a common degree $\delta\ge0$, and let
$b_i\ge0$. If $(\pi^{b_1}y_1,\ldots,\pi^{b_q}y_q,\boldsymbol t)$
is admissible, and the degrees of $\boldsymbol t$ are $\epsilon_j\ge0$,
then eventually
\begin{align*}
 &\chi(\pi^{b_1}y_1,\ldots,\pi^{b_q}y_q,\boldsymbol t;C)\\
 &\qquad=\chi(y_1,\ldots,y_q,\boldsymbol t;C)
  +\left(\sum_i b_i\right)
   \Delta_\delta^{q-1}\prod_j\Delta_{\epsilon_j}h.
\end{align*}
\end{proposition}

\begin{proof}
Apply the ordinary kernel--cokernel six-term exact sequence to
$C(-\alpha-\beta)\xrightarrow{v}C(-\alpha)\xrightarrow{u}C$,
where $\deg u=\alpha$ and $\deg v=\beta$. Additivity of the tail Euler
characteristic gives
\[
 \chi(uv,\boldsymbol t;C)_{\boldsymbol n}
 =\chi(u,\boldsymbol t;C)_{\boldsymbol n}
  +\chi(v,\boldsymbol t;C)_{\boldsymbol n-\deg u}.
\]
This retains all torsion kernels. The six modules are killed by $uv$,
so product admissibility supplies finite-length tail homology in large
degree. Peel off one degree-zero factor $u=\pi$ at a time.
Injectivity of $\pi$ identifies its Koszul factor with $C/\pi C$.
On this quotient, the Koszul terms themselves have finite length, so their
alternating lengths depend only on the degrees of the remaining sections
and give the stated backward difference. Every peeled factor contributes
that same difference. The intermediate lists are admissible because their
section ideals contain the original ideal.
\end{proof}

\paragraph{The actual lower chain.}
Return to the three-graded algebra $\mathcal R$, chart ideal $J_\varphi$,
and leading quotient $C=\mathcal R/\mathfrak q$. Let $\boldsymbol f$
be the list $F_i^2t_0,G_jt_0,E_k^2t_1,E'_lt_1,F_*^2t_2$.
The required admissibility condition is the explicit containment
\begin{equation}
 (\mathfrak m_B\mathcal R I_{\mathrm{irr}})^N
 \subset J_\varphi+(\boldsymbol f,a)
 \quad\text{for some }N,
 \label{eq:selmer-actual-admissibility}
\end{equation}
where $a\in B$ has augmentation zero and $\varphi(a)\ne0$.
The Smith-normalized negative equations give
$F_i^2t_0=\pi^{2c_i}u_i^2t_0$ in $C$ for $u_i\in A_{-1}$, by
Lemma~\ref{lem:selmer-leading-piece}. Since $a\in\mathfrak q$,
$D=(\mathcal R/J_\varphi)/a$ surjects onto $C$.
Propositions~\ref{prop:selmer-koszul-positive} and
\ref{prop:selmer-koszul-scalar}, followed by
Lemma~\ref{lem:selmer-mixed-count}, yield
\[
 \chi(\boldsymbol f;D)
 \ge\chi(\boldsymbol f;C)
 \ge\left(\sum_i2c_i\right)2^{r-1+r'}
 =2^{r+r'}\sum_i c_i.
\]
On $P=\mathcal R/J_\varphi$, multiplication by $a$ is injective:
coefficientwise application of the chart embeds $P$ in a polynomial
algebra over the domain $F[[z]]$, and $\varphi(a)\ne0$. Thus
\[
 \chi(\boldsymbol f;D)
 =\chi(a,\boldsymbol f;P)
 =\chi(\boldsymbol f,a;P)
 \le\chi(a;P/(\boldsymbol f)P).
\]
This proves the lower comparison used in
\eqref{eq:selmer-euler-retained-inputs}, conditional only on the stated
model, chart, Smith data, and
\eqref{eq:selmer-actual-admissibility}. The next proposition establishes
that containment from the stated geometry of the equation scheme. No regular sequence
or flatness of the original deformation quotient was assumed.
